\section*{\Huge Abstract}
\addcontentsline{toc}{chapter}{Abstract}

Requirements Engineering (RE) transforms high-level stakeholder intent into precise technical specifications, yet initial intent is rarely complete: it is vague, fragmented and filled with unstated assumptions that propagate into costly downstream defects. Recent LLM-based multi-agent frameworks automate RE activities by dividing the work among specialized roles, but they rely on generic stakeholder agents rather than on representations of real stakeholder groups, and they execute largely sequential workflows.

This project proposes \sysname{}, a persona-driven multi-agent framework that models requirements engineering as a dynamic, non-linear deliberation workflow. Starting from a short product idea, the framework builds a research corpus of real-world documents, mines evidence-backed stakeholder perspectives from it, and grounds the simulated stakeholder personas in that evidence. Specialized agents then interview the personas, analyze the elicited requirements, write a requirement specification and a product backlog, and validate them. Rework loops, human review checkpoints and an intent router let the process return to any stage when new information arrives, and every requirement remains traceable to the interview statements and the evidence it comes from.

\todo{Summarize the system that was built (web application, architecture) and the evaluation results in two or three sentences.}
